\title{DeepSeekMath: Pushing the Limits of Mathematical Reasoning in Open Language Models}

\begin{abstract}
Mathematical reasoning remains a major hurdle for language models due to its intricate and highly organized structure. In this work, we present DeepSeekMath~7B, developed by further pre-training DeepSeek-Coder-Base-v1.5 7B on a corpus of 120B mathematics-oriented tokens, collected from Common Crawl alongside accompanying natural language and code datasets. DeepSeekMath~7B attains a remarkable 51.7\% accuracy on the advanced MATH benchmark without depending on supplementary toolkits or ensemble approaches, bringing its performance in line with that of Gemini-Ultra and GPT-4. When self-consistency is applied over 64 generated outputs, DeepSeekMath~7B achieves a score of 60.9\% on MATH. The notable improvements in mathematical reasoning by DeepSeekMath~are underpinned by two main elements: First, a robust data curation pipeline enables comprehensive utilization of large-scale web resources. Second, we propose Group Relative Policy Optimization (GRPO), an adaptation of Proximal Policy Optimization (PPO), which strengthens mathematical reasoning capabilities while efficiently managing PPO’s memory consumption.
\end{abstract}

\section{Introduction}

Large language models (LLM) have dramatically transformed how artificial intelligence approaches mathematical reasoning, catalyzing rapid progress across quantitative reasoning benchmarks \citep{MATH} and geometric reasoning assessments \citep{trinh2024solving}. Beyond that, these models have shown considerable promise in providing support for individuals tackling sophisticated mathematical tasks \citep{tao}. Nonetheless, state-of-the-art systems such as GPT-4 \citep{gpt4} and Gemini-Ultra \citep{gemini} remain closed to the public, and currently available open-source alternatives fall well short of matching their performance.

This work presents DeepSeekMath, a specialized language model for mathematics that not only markedly surpasses other open-source models in mathematical reasoning but also comes close to matching GPT-4’s performance on a variety of academic tasks. Central to this achievement is the construction of the DeepSeekMath Corpus, an extensive, high-quality pre-training resource composed of 120B math-focused tokens. The dataset is curated from the Common Crawl (CC) by leveraging a classifier based on fastText \citep{joulin2016fasttext}. The initial classifier is trained using data from OpenWebMath \citep{openwebmath} as the positive set and incorporates a broad array of unrelated webpages as negatives. The trained classifier is then utilized to extract further positive samples from the CC; these candidates undergo a manual annotation process for greater accuracy. With this refined data, the classifier is iteratively retrained, enhancing its selection abilities. Evaluation demonstrates that the resulting large-scale dataset possesses strong quality: our DeepSeekMath-Base 7B model attains a score of 64.2\% on GSM8K \citep{gsm8k} and 36.2\% on the MATH competition benchmark \citep{MATH}, thus outperforming Minerva 540B \citep{minerva}. Furthermore, because the DeepSeekMath Corpus supports multiple languages, noticeable progress is observed in Chinese mathematical evaluation datasets \citep{wei2023cmath,agieval}. We propose that our methodology in processing mathematical data can serve as a foundation for subsequent research, with ample opportunity for further advancement.

DeepSeekMath-Base is built upon DeepSeek-Coder-Base-v1.5 7B \citep{deepseek-coder}, as our experiments indicate that leveraging a model initially trained on code yields better performance than commencing with a general-purpose LLM. Additionally, our findings reveal that mathematical pre-training not only strengthens the model’s proficiency in mathematics but also leads to improvements on general benchmarks such as MMLU \citep{mmlu} and BBH \citep{bbh}, suggesting an enhancement in broader reasoning abilities.

Following the pre-training stage, we subject DeepSeekMath-Base to mathematical instruction tuning that incorporates chain-of-thought \citep{cot}, program-of-thought \citep{pot,pal}, and tool-integrated reasoning data \citep{tora}. This process produces DeepSeekMath-Instruct 7B, which outperforms all other 7B models and even rivals larger 70B open-source instruction-tuned models.

Moreover, we present Group Relative Policy Optimization (GRPO), a variant of reinforcement learning (RL) based on Proximal Policy Optimization (PPO) \citep{schulman2017proximal}. Unlike PPO, GRPO dispenses with the critic model and instead utilizes group-based score baselines, dramatically decreasing the computational demands of training. By applying GRPO exclusively on a subset of English instruction tuning datasets, we achieve notable gains over DeepSeekMath-Instruct, in both core mathematical tasks (GSM8K: 82.9\% $\rightarrow$ 88.2\%, MATH: 46.8\% $\rightarrow$ 51.7\%) and transfer settings (e.g., CMATH: 84.6\% $\rightarrow$ 88.8\%) during RL.

We also develop a unified analytical framework that connects various approaches—Rejection Sampling Fine-Tuning (RFT) \citep{yuan2023scaling}, Direct Preference Optimization (DPO) \citep{dpo}, PPO, and GRPO—demonstrating that these can be interpreted as direct or streamlined RL methodologies. Our thorough experimentation—comparing online versus offline RL, outcome versus process-based supervision, and single-turn against iterative RL—serves to elucidate the core principles of this paradigm. Finally, we discuss the underlying mechanisms by which our RL methodology boosts instruction-tuned models and outline prospective pathways for refining RL approaches using insights from this unified perspective.

\subsection{Contributions}
We present advancements in large-scale mathematical pre-training and undertake an investigation into reinforcement learning methods.

\textbf{Math Pre-Training at Scale}

\begin{itemize}[topsep=0pt]
    \item Our study demonstrates that Common Crawl, an openly available dataset, encompasses substantial mathematical content. Through a carefully constructed data selection pipeline, we assemble the DeepSeekMath Corpus—comprising 120B tokens from math-rich web sources—which is nearly seven times greater than the math subset used by Minerva \citep{minerva}, and about nine times larger than OpenWebMath’s collection \citep{openwebmath}.

    \item The DeepSeekMath-Base 7B model, trained on this dataset, achieves a level of performance on par with Minerva 540B \citep{minerva}. This suggests that scaling the number of parameters is not the sole determinant of strong mathematical reasoning; pre-training a relatively smaller model on sufficiently high-quality data can yield impressive results.

    \item Our experimental results in math model training indicate that preceding math pre-training with code pre-training enhances problem-solving capabilities—both in the presence and absence of tools. These findings lend support to the hypothesis that code pre-training boosts reasoning, specifically in mathematical domains.

    \item Despite its popularity, the practice of incorporating arXiv paper data during training—common in recent mathematical language models—yields no significant gains across the mathematical evaluation benchmarks considered in our work.
\end{itemize}

\textbf{Exploration and Analysis of Reinforcement Learning}

\begin{itemize}[topsep=0pt]
    \item This work proposes Group Relative Policy Optimization (GRPO), a reinforcement learning method that emphasizes both efficiency and effectiveness. Distinct from Proximal Policy Optimization (PPO), GRPO eliminates the need for a critic network, deriving the baseline from aggregate group scores instead. This modification leads to a substantial reduction in computational demands during training.
    \item Our experiments show that GRPO notably improves the performance of DeepSeekMath-Instruct, our instruction-tuned model, using only the data from instruction-tuning. Additionally, we identify performance gains on out-of-domain tasks during reinforcement learning, indicating a broader generalization benefit.
    \item We articulate a cohesive conceptual framework that encompasses various techniques, such as RFT, DPO, PPO, and GRPO. To dissect the key contributors to this framework, we carry out extensive empirical investigations—covering comparisons between online and offline learning, supervision over outcomes versus processes, single-turn versus iterative RL, among others.
    \item Leveraging insights from this unified perspective, we analyze factors responsible for effective reinforcement learning, and propose several promising avenues for advancing the reinforcement learning capabilities of large language models (LLMs).
\end{itemize}

\subsection{Summary of Evaluations and Metrics}

\begin{itemize}[topsep=0pt]
    \item \textbf{English and Chinese Mathematical Reasoning}: We conduct thorough evaluations of our models across both English and Chinese benchmarks, spanning a range of mathematical tasks from elementary to collegiate complexity. The English datasets encompass GSM8K \citep{gsm8k}, MATH \citep{MATH}, SAT \citep{llemma}, OCW Courses \citep{minerva}, and MMLU-STEM \citep{mmlu}, while the Chinese tasks utilize MGSM-zh \citep{mgsm}, CMATH \citep{wei2023cmath}, Gaokao-MathCloze \citep{agieval}, and Gaokao-MathQA \citep{agieval}. Our evaluation protocol assesses both the models' capability to independently compose complete solutions in natural language, as well as their proficiency in solving problems programmatically using Python.

    On the English mathematical benchmarks, DeepSeekMath-Base demonstrates performance comparable to the proprietary Minerva 540B \citep{minerva}, and outperforms all publicly available base models—including Mistral 7B \citep{mistral} and Llemma-34B \citep{llemma}—regardless of prior math-specific pre-training, and often does so with a considerable lead. Particularly on the Chinese benchmarks, DeepSeekMath-Base achieves top results, attributed to our divergence from the English-only pre-training focus of earlier studies \citep{minerva,llemma} by incorporating superior multilingual data. When enhanced with instruction tuning and reinforcement learning, the variants DeepSeekMath-Instruct and DeepSeekMath-RL attain state-of-the-art results, becoming the first open-source models to surpass 50\% accuracy on the competitive MATH benchmark.

\item \textbf{Formal Mathematics}: DeepSeekMath-Base is assessed on the informal-to-formal theorem proving benchmark presented in \citep{dsp_proof}, utilizing miniF2F \citep{minif2f} and Isabelle \citep{isabelle} as the chosen proof assistant. In this setting, DeepSeekMath-Base achieves impressive results in few-shot autoformalization tasks.

\item \textbf{Natural Language Understanding, Reasoning, and Code}: In order to thoroughly evaluate the general proficiency of models in understanding, reasoning, and coding, we benchmark DeepSeekMath-Base on the Massive Multitask Language Understanding (MMLU) benchmark \citep{mmlu}, which includes 57 diverse multiple-choice tasks, as well as BIG-Bench Hard (BBH) \citep{bbh}, a collection of 23 tasks primarily designed to test complex multi-step reasoning. Additionally, we evaluate performance on code generation using HumanEval \citep{codex} and MBPP \citep{mbpp}, two established standards in the field. The results show that pre-training on mathematical data enhances both reasoning and language understanding capabilities.

\section{Math Pre-Training}

\subsection{Data Collection and Decontamination}
Here, we describe our methodology for assembling the DeepSeekMath Corpus based on data from Common Crawl. As illustrated in Figure \ref{fig:data_collect}, our workflow adopts an iterative strategy to systematically construct a large-scale mathematical dataset from Common Crawl, beginning with a seed corpus such as a curated collection of math-related material. This procedure is not limited to mathematics and can be generalized to other fields, including programming.

To initiate data collection, we select OpenWebMath \citep{openwebmath}, which provides a repository of curated mathematical web content, to serve as the foundational seed corpus. We utilize this initial data to train a fastText model \citep{joulin2016fasttext} with the objective of retrieving additional web pages that share similar mathematical characteristics to those found in OpenWebMath. Specifically, we designate 500,000 instances from the seed dataset as positive samples, while another 500,000 pages from Common Crawl serve as negative counterparts. Training is conducted using an open-source implementation\footnote{\url{https://fasttext.cc}}, with the vector size set to 256, a learning rate of 0.1, a maximum word n-gram length of 3, a minimum word occurrence of 3, and 3 epochs. To address data redundancy, we first apply both exact and near-duplicate filtering based on URLs, reducing the Common Crawl dataset to 40B HTML pages. The trained fastText model is then used to extract mathematical content from this deduplicated subset. Collected pages are subsequently ranked by their fastText-predicted scores, ensuring only the highest quality pages are retained. The effectiveness of this filtering is assessed by pre-training on the top 40B, 80B, 120B, and 160B tokens, with the first iteration focused on the top 40B tokens for downstream use.

Following the initial round of data collection, a substantial number of mathematical web pages remain uncollected. This is largely due to the limited diversity present in the set of positive examples used to train the fastText model. To address this shortcoming and improve model performance, we incorporate additional mathematical web domains into the seed corpus. Concretely, we begin by dividing the entirety of Common Crawl into separate domains, where each domain is defined as the set of web pages that share an identical base URL. For every domain, we determine the proportion of web pages retrieved during the first collection iteration. Any domain from which more than 10\% of its pages were retrieved is designated as math-related (e.g., \url{mathoverflow.net}). 

We then proceed to manually label URLs associated with mathematical content within these math-related domains (such as \url{mathoverflow.net/questions}). Previously uncollected pages from these URLs are incorporated into the seed corpus, broadening the range of positive training examples. This expanded dataset serves to enhance the performance of the fastText model, thereby improving its ability to identify and retrieve mathematical content in subsequent collection rounds. After repeating this process over four iterations, the collected dataset comprises 35.5M mathematical web pages, amounting to 120B tokens. By the fourth iteration, approximately 98\% of the data had already been captured by the third iteration, prompting the termination of additional data collection.

In order to prevent contamination of benchmark datasets, we adopt the methodology outlined in \cite{deepseek-coder}, filtering out any web pages that contain questions or answers appearing in English mathematical benchmarks—such as GSM8K \citep{gsm8k} and MATH \citep{MATH}—as well as in Chinese benchmarks like CMATH \citep{wei2023cmath} and AGIEval \citep{agieval}. The specific filtering protocol excludes any text span containing a 10-gram that exactly matches any substring found in the evaluation benchmarks from our mathematical training dataset. For benchmark texts comprising fewer than 10 but at least 3 grams, an exact match criterion is similarly employed to exclude contaminated web pages.

\subsection{Validation of the DeepSeekMath~Corpus Quality}

We conduct pre-training studies to compare the DeepSeekMath~Corpus against recently released math-focused training corpora:

\begin{itemize}[topsep=0pt]
    \item \textbf{MathPile} \citep{mathpile}: A composite corpus of 8.9B tokens drawing from various sources, including textbooks, Wikipedia, ProofWiki, CommonCrawl, StackExchange, and arXiv, where arXiv content accounts for over 85\%;
    \item \textbf{OpenWebMath} \citep{openwebmath}: Data extracted from CommonCrawl and curated to select mathematical material, totaling 13.6B tokens;
    \item \textbf{Proof-Pile-2} \citep{llemma}: This corpus combines OpenWebMath, AlgebraicStack (10.3B tokens from mathematical code), and arXiv (28.0B tokens). In line with \cite{llemma}, we experiment with the arXiv:Web:Code composition set to a ratio of 2:4:1.
\end{itemize}

\subsubsection{Training Setting}
\label{sec:quality-policy}
We fine-tune a general-purpose language model containing 1.3B parameters—architecturally consistent with the DeepSeek LLMs \citep{deepseek-llm} and henceforth referenced as DeepSeek-LLM 1.3B—on mathematical data. Each individual mathematical corpus serves as the basis for a separate model, with each being trained for a total of 150B tokens. For all training, we utilize the efficient HAI-LLM framework \citep{haillm}. Following protocols established in prior DeepSeek LLM work, optimization is conducted using AdamW \citep{adamW}, setting $\beta_1=0.9$, $\beta_2=0.95$, and a $\mathrm{weight\_decay}=0.1$. The learning rate scheduling incorporates a multi-phase strategy: the rate is linearly ramped up to its maximum over 2{,}000 warmup steps, then reduced to 31.6\% of its peak value by 80\% of training, and finally decays to 10.0\% of the maximum by the final 10\% of training. We select 5.3e-4 as the highest learning rate, adopting a batch size equivalent to 4M tokens and a 4K token context window.

\subsubsection{Evaluation Results}

\textbf{DeepSeekMath~Corpus excels in quality, breadth of multilingual mathematical materials, and scale.}

\begin{itemize}[topsep=0pt]
    \item \textbf{High-quality}:
    We assess downstream model capabilities on 8 mathematical benchmarks leveraging few-shot chain-of-thought prompting \cite{cot}. According to Table \ref{tab:corpora_comparison}, the model trained with the DeepSeekMath~Corpus distinctly outperforms other models. Additionally, Figure \ref{fig:corpora_comparisons} reveals that, at the 50B token mark (corresponding to a full pass over Proof-Pile-2), the DeepSeekMath-trained model surpasses its counterpart trained on Proof-Pile-2. This finding underscores the superior average quality of the DeepSeekMath Corpus.
    \item \textbf{Multilingual}:
    DeepSeekMath~Corpus includes content in a variety of languages, with English and Chinese being especially prominent. Table \ref{tab:corpora_comparison} demonstrates that DeepSeekMath~Corpus boosts mathematical reasoning for both English and Chinese tasks. Conversely, existing corpora, mainly focusing on English, show only marginal gains and can even adversely affect Chinese mathematical reasoning capabilities.
    \item \textbf{Large-scale}:
    In terms of size, DeepSeekMath~Corpus greatly exceeds current mathematical corpora. As illustrated in Figure \ref{fig:corpora_comparisons}, DeepSeek-LLM 1.3B, when fine-tuned on DeepSeekMath~Corpus, exhibits a more rapid ascent in learning as well as more sustained progress. In comparison, the baseline datasets are considerably smaller, necessitating repeated epochs during training, which results in the associated model performance quickly saturating.
\end{itemize}

\subsection{Training and Evaluating DeepSeekMath-Base 7B}

This section presents DeepSeekMath-Base 7B, a foundational model designed to demonstrate advanced reasoning, particularly within mathematical domains. The model initialization employs DeepSeek-Coder-Base-v1.5 7B \citep{deepseek-coder}, subsequently undergoing training on a dataset totaling 500B tokens. The composition of this training data is allocated as follows: DeepSeekMath Corpus accounts for 56\%, AlgebraicStack for 4\%, arXiv for 10\%, Github code constitutes 20\%, while the final 10\% consists of natural language data sourced from Common Crawl in both English and Chinese. Training follows the general protocol outlined in Section \ref{sec:quality-policy}, with the primary modifications being an increased peak learning rate of 4.2e-4 and the application of a batch size comprising 10M tokens.

We systematically evaluate DeepSeekMath-Base 7B's mathematical competence by assessing its aptitude for independently constructing mathematical solutions, employing auxiliary tools for problem-solving, and executing formal theorem proving tasks. To complement this mathematical evaluation, we also investigate broader capabilities of the model, such as proficiency in natural language comprehension, logical reasoning, and coding.

\paragraph{Mathematical Problem Solving with Step-by-Step Reasoning}

To measure the effectiveness of DeepSeekMath-Base in mathematical problem-solving, we utilize few-shot chain-of-thought prompting \citep{cot} across a suite of eight benchmarks provided in both English and Chinese. These benchmarks include quantitative reasoning assessments (such as GSM8K \citep{gsm8k}, MATH \citep{MATH}, and CMATH \citep{wei2023cmath}), as well as multiple-choice formats (including MMLU-STEM \citep{mmlu} and Gaokao-MathQA \citep{agieval}), thereby spanning a wide spectrum of mathematical topics, from foundational arithmetic to advanced collegiate material.

As detailed in Table \ref{tab:base_math_cot}, DeepSeekMath-Base 7B consistently achieves the highest scores across all eight benchmark tasks among publicly available base models—this includes prominent models such as the broadly adopted Mistral 7B \citep{mistral} and the newly introduced Llemma 34B \citep{llemma}, which received mathematical training via Proof-Pile-2 \citep{llemma}. Of particular interest, DeepSeekMath-Base achieves more than a 10\% absolute improvement over other open-source base models on the MATH competition benchmark, even outperforming Minerva 540B \citep{minerva}—a substantially larger proprietary model (77 times greater in size) built on the PaLM architecture \citep{palm} and further pre-trained on math-centric data.

\paragraph{Mathematical Problem Solving with Tool Use} 

To assess mathematical reasoning facilitated by external tools, we apply few-shot program-of-thought prompting \citep{pot,pal} on GSM8K and MATH. For each problem, models are prompted to generate Python programs that can employ modules like \textit{math} or \textit{sympy} for complex calculations. The numerical result returned by executing the generated program is then regarded as the solution. According to the findings in Table \ref{tab:base_math_pot_proof}, DeepSeekMath-Base 7B attains higher accuracy than the previous state-of-the-art Llemma 34B.

\paragraph{Formal Mathematics}

Automating formal proof processes is increasingly recognized for its capacity to guarantee correctness in mathematical arguments and to streamline proof construction. We benchmark DeepSeekMath-Base 7B using the informal-to-formal proving task described in \citep{dsp_proof}, which asks for the production of formal proofs based on an informal assertion, its formal version, and a corresponding informal proof. The evaluation utilizes miniF2F \citep{minif2f}, a collection representing formal Olympiad-style mathematics challenges, and involves generating Isabelle-formatted formal proofs through few-shot prompting. In line with \cite{dsp_proof}, our method has the model produce proof outlines, subsequently invoking Sledgehammer \citep{sledgehammer}—an automated theorem prover—to complete any proof gaps. Table \ref{tab:base_math_pot_proof} illustrates the strong performance of DeepSeekMath-Base 7B in automating the formalization of proofs.

\paragraph{Natural Language Understanding, Reasoning, and Code}

For a broader assessment, we evaluate the model’s proficiency in language comprehension on MMLU \citep{mmlu}, complex reasoning on BBH \citep{bbh}, and coding on HumanEval \citep{codex} and MBPP \citep{mbpp}. As shown in Table \ref{tab:understanding_reasoning_code}, DeepSeekMath-Base 7B delivers notable improvements over its predecessor, DeepSeek-Coder-Base-v1.5 \citep{deepseek-coder}, in both language understanding (MMLU) and reasoning (BBH), demonstrating the beneficial effects of mathematical pre-training in these areas. Furthermore, owing to the continuous exposure to code tokens during training, DeepSeekMath-Base 7B preserves the coding task performance established by DeepSeek-Coder-Base-v1.5 on HumanEval and MBPP. Overall, the model decisively surpasses Mistral 7B \citep{mistral} across the evaluated reasoning and coding benchmarks.

\section{Supervised Fine-Tuning}

\subsection{SFT Data Curation}

To assemble a mathematical instruction-tuning dataset, we gather English and Chinese problem sets across a range of mathematical domains and difficulty levels. Each problem is aligned with its solution in various reasoning modalities, including chain-of-thought (CoT) \citep{cot}, program-of-thought (PoT) \citep{pot,pal}, and solutions utilizing external tools \citep{tora}. The compiled training set consists of 776,000 instances.

\begin{itemize}[topsep=0pt]
    \item \textbf{English mathematical datasets}: We enhance GSM8K and MATH problem sets by annotating them with tool-based solutions. Additionally, we incorporate portions of MathInstruct \citep{MathInstruct}, as well as the training partition from Lila-OOD \citep{lila}, featuring solutions produced via CoT or PoT methods. This English dataset spans a breadth of mathematical specialties, including algebra, calculus, number theory, probability, and geometry.
    \item \textbf{Chinese mathematical datasets}: Our Chinese corpus includes K-12 level math problems encompassing 76 sub-domains such as linear equations, each paired with explanations formatted in both CoT and tool-enhanced reasoning styles.
\end{itemize}

\subsection{Training and Evaluating DeepSeekMath-Instruct 7B}

Here, we describe DeepSeekMath-Instruct 7B, which is mathematically instruction-tuned using DeepSeekMath-Base as the foundation. Training data are concatenated in a random order, not exceeding a sequence length of 4,000 tokens per batch. The model undergoes 500 optimization steps, employing a batch size of 256 and a fixed learning rate of $5\times10^{-5}$.

We assess the mathematical capabilities of the models, both in the absence and presence of tool assistance, using four quantitative reasoning benchmarks presented in English and Chinese. For comparison, we evaluate against contemporary top-performing models:

\begin{itemize}[topsep=0pt]
    \item \textbf{Closed-source models} encompass: (1) models from the GPT series, including GPT-4 \citep{gpt4} and its Code Interpreter variant~\footnote{\url{https://openai.com/blog/chatgpt-plugins\#\#code-interpreter}}, (2) Gemini Ultra and Pro \citep{gemini}, (3) Inflection-2 \citep{inflection-2}, (4) Grok-1~\footnote{\url{https://x.ai/model-card}}, and (5) recent entrants from Chinese developers such as Baichuan-3~\footnote{\url{https://www.baichuan-ai.com}}, and (6) the newest GLM-4~\footnote{\url{https://open.bigmodel.cn/dev/api\#glm-4}} from the GLM model line \citep{glm}.
\end{itemize}

These models are intended for broad applicability, with most having undergone various alignment processes.

\item \textbf{Open-source models} consist of: general-purpose models such as (1) DeepSeek-LLM-Chat 67B \citep{deepseek-llm}, (2) Qwen 72B \citep{qwen}, (3) SeaLLM-v2 7B \citep{seallm}, and (4) ChatGLM3 6B \citep{chatglm3}. Additionally, several models incorporate mathematics-oriented enhancements, including: (5) InternLM2-Math 20B~\footnote{\url{https://github.com/InternLM/InternLM-Math}}, which extends InternLM2 by conducting specialized math training and subsequent instruction fine-tuning; (6) Math-Shepherd-Mistral 7B, which utilizes PPO training \citep{schulman2017proximal} on Mistral 7B \citep{mistral} aided by a process-supervised reward model; (7) the WizardMath series \citep{wizardmath}, which targets improved mathematical reasoning for Mistral 7B and Llama-2 70B \citep{llama2} using evolve-instruct (an instruction tuning method leveraging AI-generated instructions) and PPO optimization, primarily trained on problems from GSM8K and MATH; (8) MetaMath 70B \citep{metamath}, which adapts Llama-2 70B with an augmented GSM8K and MATH dataset; (9) ToRA 34B \cite{tora}, which tailors CodeLlama 34B for tool-augmented mathematical reasoning through fine-tuning; (10) MAmmoTH 70B \citep{MathInstruct}, which involves instruction tuning Llama-2 70B using MathInstruct.

Referencing Table \ref{tab:sft_rl_math}, when tool use is not permitted in the evaluation, DeepSeekMath-Instruct 7B achieves notable step-by-step reasoning capability. Specifically, on the advanced MATH benchmark, it exceeds the accuracy of all open-source and most proprietary competitors (such as Inflection-2 and Gemini Pro) by a minimum margin of 9\% absolute—this includes models significantly larger in scale (e.g., Qwen 72B) or those benefiting from intensive math-specific reinforcement learning (e.g., WizardMath-v1.1 7B). While DeepSeekMath-Instruct is on par with the proprietary Chinese systems GLM-4 and Baichuan-3 in MATH performance, it remains behind GPT-4 and Gemini Ultra.

In scenarios permitting both natural language and code-based tool reasoning, DeepSeekMath-Instruct 7B achieves an accuracy close to 60\% on the MATH dataset, outperforming all open-source alternatives. Across other tasks, it matches or exceeds the performance of DeepSeek-LLM-Chat 67B, the previously leading open-source model of ten times greater parameter count.

\section{Reinforcement Learning}

\subsection{Group Relative Policy Optimization} Reinforcement learning (RL) has consistently demonstrated its capacity to further strengthen LLMs' mathematical reasoning after supervised fine-tuning (SFT) \citep{wang2023math, wizardmath}. In this context, we present our streamlined yet powerful RL method: Group Relative Policy Optimization (GRPO).

\subsubsection{From PPO to GRPO} Proximal Policy Optimization (PPO) \citep{schulman2017proximal} represents a widely adopted actor-critic reinforcement learning strategy for LLM fine-tuning \citep{ouyang2022training}. This algorithm refines LLM policies by maximizing the following surrogate objective:

\begin{equation}
\footnotesize
    \mathcal{J}_{PPO}(\theta) = \mathbb{E}{[q \sim P(Q), o \sim \pi_{\theta_{old}}(O|q)]} \frac{1}{|o|} \sum_{t=1}^{|o|} \min \left[ \frac{\pi_\theta(o_{t} | q, o_{<t})}{\pi_{\theta_{old}}(o_{t} | q, o_{<t})} A_{t}, \text{clip} \left( \frac{\pi_\theta(o_{t} | q, o_{<t})}{\pi_{\theta_{old}}(o_{t} | q, o_{<t})}, 1 - \varepsilon, 1 + \varepsilon \right)  A_{t} \right] ,
\end{equation}

Here, $\pi_{\theta}$ and $\pi_{\theta_{old}}$ denote the current and preceding policy models, respectively, while $q$ and $o$ represent questions and responses sampled from the question corpus and the older policy $\pi_{\theta_{old}}$. The parameter $\varepsilon$ serves as a clipping hyper-parameter in PPO, designed to enhance the robustness of training. $A_t$ signifies the advantage estimate, obtained using Generalized Advantage Estimation (GAE) \citep{gae}, based on reward signals $\{r_{\ge t}\}$ and a learned value function $V_{\psi}$. Consequently, PPO necessitates parallel training of a value function alongside the policy, and to control overfitting to the reward model, the typical practice is to include a per-token KL-divergence penalty (relative to a reference model) in the reward at each token, following \citep{ouyang2022training}:

\begin{equation}
    r_{t} = r_\varphi(q, o_{\le t}) - \beta \log\frac{\pi_{\theta}(o_{t}|q, o_{<t})}{\pi_{ref}(o_{t}|q, o_{<t})},
\label{eq:PPO-reward}
\end{equation}

In this formula, $r_\varphi$ stands for the reward model, $\pi_{ref}$ indicates the reference model—often set to the original SFT model—and $\beta$ is the scaling factor for the KL regularization term.

Because the value network employed in PPO is usually as large as the policy network, its inclusion results in considerable increases in both memory consumption and computational requirements. Moreover, during reinforcement learning, this value network functions as a baseline to decrease variance in the advantage estimation. However, in the large language model setting, it is common for the reward model to assign non-zero reward only to the final token, which poses challenges in producing accurate token-level value estimates. To address these issues, we introduce, as illustrated in Figure \ref{fig:grpo}, a new method named Group Relative Policy Optimization (GRPO). GRPO eliminates the need for learning an extra value function as in PPO; instead, it takes the mean reward across several samples generated for a single question as a baseline. Concretely, for each prompt $q$, a set of outputs $\{o_1, o_2, \cdots, o_G\}$ is generated using the previous policy $\pi_{\theta_{old}}$, and the policy is then updated to maximize:

\begin{equation}
\footnotesize
\begin{split}
    \mathcal{J}_{GRPO}(\theta) &= \mathbb{E}{[q \sim P(Q), \{o_i\}_{i=1}^G \sim \pi_{\theta_{old}}(O|q)]}  \\
    & \frac{1}{G}\sum_{i=1}^G\frac{1}{|o_i|} \sum_{t=1}^{|o_i|} \left\{ \min \left[ \frac{\pi_\theta(o_{i,t} | q, o_{i,<t})}{\pi_{\theta_{old}}(o_{i,t} | q, o_{i,<t})} \hat{A}_{i,t}, \text{clip} \left( \frac{\pi_\theta(o_{i,t} | q, o_{i,<t})}{\pi_{\theta_{old}}(o_{i,t} | q, o_{i,<t})}, 1 - \varepsilon, 1 + \varepsilon \right)  \hat{A}_{i,t} \right] - \beta \mathbb{D}_{KL}\left[\pi_{\theta} || \pi_{ref}\right]\right\} ,
\end{split}
\label{eq:GRPO-obj}
\end{equation}

Here, $\varepsilon$ and $\beta$ remain as hyper-parameters, while $\hat{A}_{i,t}$ corresponds to the advantage, which is computed exclusively based on relative scores within each sampled group (elaborated in subsequent sections). The design of GRPO, which computes advantages in a group-relative fashion, aligns with the nature of reward models, since such models are generally trained using comparisons between outputs given the same prompt. Another important aspect is that GRPO applies the KL penalty term directly at the loss level rather than incorporating it into the reward calculation, thereby simplifying the derivation of $\hat{A}_{i,t}$. Moreover, instead of the KL penalty defined in (\ref{eq:PPO-reward}), GRPO employs an unbiased estimator for the KL-divergence as suggested by \citep

\begin{algorithm}[t]
  \small
  \caption{Iterative Group Relative Policy Optimization}
  \textbf{Input} initial policy model $\pi_{\theta_{\text{init}}}$; reward models $r_\varphi$; task prompts $\mathcal{D}$; 
  hyperparameters $\varepsilon$, $\beta$, $\mu$
  \begin{algorithmic}[1]
    \State Initialize the policy model: $\pi_\theta \leftarrow \pi_{\theta_{\text{init}}}$
    \For{iteration = 1, \dots, I}
       \State Set reference model $\pi_{ref} \leftarrow \pi_{\theta}$
      \For{step = 1, \dots, M}
      \State Draw a batch $\mathcal{D}_b$ from $\mathcal{D}$
      \State Update previous policy: $\pi_{\theta_{old}} \leftarrow \pi_{\theta}$ 
      \State For each question $q \in \mathcal{D}_b$, generate $G$ samples $\{o_i\}_{i=1}^G \sim \pi_{\theta_{old}} (\cdot \mid q) $
      \State Evaluate each sampled output $o_i$ with the reward model $r_\varphi$ to obtain rewards $\{r_i\}_{i=1}^{G}$
      \State Determine $\hat{A}_{i,t}$ for the $t$-th token of $o_i$ using group relative advantage estimation
      \For{GRPO iteration = 1, \dots, $\mu$}
        \State Optimize the policy model $\pi_\theta$ by maximizing the GRPO objective (Equation \ref{eq:GC-GRPO})
      \EndFor
    \EndFor \State Continuously train $r_\varphi$ using a replay-based mechanism for updates. 
    \EndFor 
  \end{algorithmic}
  \textbf{Output} $\pi_\theta$
  \label{alg:iter-grpo}
\end{algorithm}

\subsubsection{Outcome Supervision RL with GRPO} In this setup, given a question $q$, a set of $G$ candidate outputs $\{o_1, o_2, \cdots, o_G\}$ is sampled using the previous policy $\pi_{\theta_{old}}$. Each output is assessed by a reward model, resulting in a set of rewards $\mathbf{r}=\{r_1, r_2, \cdots, r_G\}$. These rewards are standardized by subtracting the mean and dividing by the standard deviation of the group. For outcome supervision, the standardized reward is provided at the end of each output $o_i$, and this value is uniformly assigned as the advantage $\hat{A}_{i, t}$ to all tokens within the output, that is, $\hat{A}_{i, t} = \widetilde{r}_i = \frac{r_i- {\rm mean}(\mathbf{r})}{{\rm std}(\mathbf{r})}$. Policy optimization proceeds by maximizing the target specified in equation (\ref{eq:GRPO-obj}).

\subsubsection{Process Supervision RL with GRPO} Providing only terminal rewards, as in outcome supervision, can be insufficient or inefficient for complex mathematical problem solving. Following \cite{wang2023math}, we incorporate process supervision, which assigns rewards at the end of each intermediate reasoning step. Specifically, for a question $q$ and a set of $G$ sampled responses $\{o_1, o_2, \cdots, o_G\}$, a process reward model computes a reward at each step, forming reward sequences: $\mathbf{R} = \{ \{r_1^{index(1)},\cdots,r_1^{index(K_1)}\}, \cdots,  \{r_G^{index(1)},\cdots,r_G^{index(K_G)}\} \}$, where $index(j)$ denotes the token index marking the end of the $j$-th step, and $K_i$ is the total number of steps in $o_i$. Each step's reward is normalized using the group’s mean and standard deviation: $\widetilde{r}_i^{index(j)} = \frac{r_i^{index(j)} - {\rm mean(\mathbf{R})}}{{\rm std(\mathbf{R})}}$.
Then, for process supervision, the advantage for token $t$ is the cumulative sum of normalized rewards from the steps following $t$, i.e., $\hat{A}_{i, t} = \sum_{index(j) \ge t} \widetilde{r}_i

\subsubsection{Iterative RL with GRPO} During reinforcement learning, the previous reward model might not adequately guide the evolving policy model. To address this, we investigate an iterative RL strategy incorporating GRPO. As depicted in Algorithm \ref{alg:iter-grpo}, the iterative GRPO framework involves producing updated training datasets for the reward model from samples drawn by the policy model. The reward model is then further refined using a replay strategy that includes 10\% of its historical samples. Subsequently, the policy model itself becomes the reference model and undergoes continued optimization, now leveraging the updated reward model.

\subsection{Training and Evaluating DeepSeekMath-RL}

Our RL experiments utilize DeepSeekMath-Instruct 7B as the foundation. The reinforcement learning phase draws on chain-of-thought-formatted GSM8K and MATH questions from the SFT data, totaling approximately 144K items. To specifically examine RL's influence on benchmarks absent during RL training, all other SFT items are omitted. The reward model’s training set is constructed in accordance with \citep{wang2023math}. The initial reward model is fine-tuned from DeepSeekMath-Base 7B using a learning rate of 2e-5. For GRPO, the policy model employs a learning rate of 1e-6, and a KL coefficient of 0.04. Each question is used to sample $64$ completions. Maximum sequence length is set to 1024, with a batch size of 1024. The policy model is updated just once after each exploration step. Evaluation of DeepSeekMath-RL 7B follows the same benchmarks as DeepSeekMath-Instruct 7B, designating GSM8K and MATH with chain-of-thought reasoning as in-domain, while treating all other benchmarks as out-of-domain tasks.

Table \ref{tab:sft_rl_math} presents a comparative analysis of both open- and closed-source models utilizing chain-of-thought and tool-augmented reasoning approaches on English and Chinese benchmark datasets. Our observations reveal the following: 1) With chain-of-thought prompting, DeepSeekMath-RL 7B achieves 88.2\% accuracy on GSM8K and 51.7\% on MATH, outperforming every open-source model in the 7B to 70B parameter range and exceeding the capabilities of most closed-source models as well. 2) Notably, DeepSeekMath-RL 7B is exclusively trained on chain-of-thought-format instruction data from GSM8K and MATH, and its initial checkpoint is DeepSeekMath-Instruct 7B. Although the training data is narrowly focused, DeepSeekMath-RL 7B consistently surpasses DeepSeekMath-Instruct 7B across evaluation metrics, thereby illustrating the pronounced benefit of reinforcement learning.

\section{Discussion}
Here, we elaborate on the insights gathered during our pre-training and reinforcement learning experiments.

\subsection{Lessons Learnt in Pre-Training}

We begin by detailing the insights derived from our pre-training process. Unless otherwise indicated, the configurations follow the standards described in Section \ref{sec:quality-policy}. Importantly, when mentioning the DeepSeekMath Corpus within this discussion, we refer specifically to the 89B-token dataset generated during the second phase of data collection.

\subsubsection{Code Training Benefits Mathematical Reasoning}

There exists a widely held but largely untested assumption that code pre-training can bolster reasoning abilities. Our investigations partially substantiate this in the context of mathematical reasoning: pre-training on code appears to enhance a model’s mathematical reasoning skills, both with and without the aid of external tools.

To explore the influence of code training on mathematical reasoning capabilities, we considered two experimental setups: one involving a two-stage training process and another using a one-stage approach:

\textbf{Two-Stage Training}

\begin{itemize}[topsep=0pt]
    \item \textbf{Code Training for 400B Tokens $\rightarrow$ Math Training for 150B Tokens}: DeepSeek-LLM 1.3B is first pretrained on 400B tokens of code data, and then further trained on 150B math tokens;
    \item \textbf{General Training for 400B Tokens $\rightarrow$ Math Training for 150B Tokens}: For comparative analysis, the model is initially exposed to 400B general tokens (drawn from an extensive general-purpose corpus assembled by DeepSeek-AI) instead of code tokens before proceeding with math-specific training, to evaluate the relative efficacy of code versus general-domain pretraining on mathematical reasoning.
\end{itemize}

\textbf{One-Stage Training}

\begin{itemize}[topsep=0pt]
    \item \textbf{Math Training for 150B Tokens}: Here, DeepSeek-LLM 1.3B undergoes training exclusively on 150B math tokens;
    \item \textbf{Training on a mixture of 400B Code Tokens and 150B Math Tokens}: Since subsequent math-focused training after code pretraining often diminishes performance on coding tasks, we assess whether concurrently presenting the model with both code and math data can yield enhancements in mathematical reasoning while also mitigating the challenge of catastrophic forgetting.
\end{itemize}

\paragraph{Results}
The outcomes under different training schemes are summarized in Table \ref{tab:code-math} and Table \ref{tab:code-math-general-eval}.

Incorporating code pretraining yields substantial improvements in program-based mathematical reasoning across both two-stage and mixed-stage regimes. Specifically, Table \ref{tab:code-math} reveals that pretraining on code alone, as the first stage, greatly elevates the model’s proficiency on GSM8K and MATH benchmarks involving Python reasoning; further advancements are observed after the math-specific phase. Moreover, conducting simultaneous training on code and math tokens in a single stage efficiently reduces catastrophic forgetting—an issue prominent with sequential training—while additionally supporting robust performance in both code-related tasks (Table \ref{tab:code-math-general-eval}) and program-assisted mathematical problem solving (Table \ref{tab:code-math}).

Code-centric pretraining also enhances non-tool-based mathematical reasoning, as evidenced by two-stage training: initial code exposure yields moderate gains and improves the effectiveness of subsequent math training, culminating in peak performance. In contrast, fusing code and math data in a single-stage regime appears detrimental to pure mathematical reasoning, potentially due to DeepSeek-LLM 1.3B’s limited parameter capacity inhibiting the thorough absorption of both data modalities when learned in tandem.

\subsubsection{ArXiv Papers Appear Ineffectual for Enhancing Mathematical Reasoning}

While arXiv papers are often incorporated into mathematical pre-training datasets \citep{minerva,gpt-f,llemma,mathpile}, there has been little rigorous investigation into their actual influence on mathematical reasoning performance. Contrary to perhaps common intuition, our empirical results suggest that pre-training on arXiv papers does not substantially benefit mathematical reasoning capabilities. To evaluate this, we considered models of various capacities, namely DeepSeek-LLM 1.3B and DeepSeek-Coder-Base-v1.5 7B \citep{deepseek-coder}, pre-trained on differently curated arXiv collections:

\begin{itemize}[topsep=0pt]
    \item \textbf{MathPile} \citep{mathpile}: an 8.9B-token dataset assembled via cleaning and filtering heuristics, with over 85\% comprised of arXiv scientific papers;
    \item \textbf{ArXiv-RedPajama} \citep{redpajama}: a 28.0B-token dump of all arXiv LaTeX submissions, purged of preambles, macros, comments, and bibliographies.
\end{itemize}

Our approach involved training DeepSeek-LLM 1.3B for 150B tokens and DeepSeek-Coder-Base-v1.5 7B for 40B tokens on each arXiv-based dataset independently. Across these experiments, neither model showed meaningful gains in mathematical reasoning or, in certain scenarios, even exhibited worsened performance on a suite of mathematical benchmarks of varying difficulty. The benchmarks covered include quantitative reasoning tasks such as GSM8K and MATH (see Table \ref{tab:arxiv-cot}), multiple-choice STEM evaluations like MMLU-STEM (Table \ref{tab:arxiv-cot}), and formal mathematics problems as found in miniF2F (Table \ref{tab:arxiv_atp}).

Nonetheless, these findings must be interpreted cautiously due to several caveats. Our study does not address:

\begin{itemize}[topsep=0pt]
    \item The influence of arXiv-derived data on math tasks beyond those assessed here, such as transforming formal mathematical content into its informal counterpart;
    \item Interactions between arXiv data and other forms of training data;
    \item Potential positive effects that may emerge at greater model scales.
\end{itemize}

Accordingly, these aspects warrant further investigation, which we defer to subsequent research.

\subsection{Insights into Reinforcement Learning}
\subsubsection{Toward a Unified Training Paradigm} Here, we introduce a unified perspective for comparing various training methods—namely SFT, RFT, DPO, PPO, and GRPO—and conduct empirical studies on aspects of this framework. Generally, the gradient of a training objective with respect to parameters $\theta$ can be formulated as:

\begin{equation}
    \nabla_{\theta}\mathcal{J}_{\textcolor{red}{\mathcal{A}}}(\theta) = \mathbb{E}[\underbrace{(q,o) \sim \textcolor{red}{\mathcal{D}}}_{Data \ Source}]\left( \frac{1}{|o|} \sum_{t=1}^{|o|}  \underbrace{GC_{{\mathcal{A}}}(q, o, t, \textcolor{red}{\pi_{{rf}}})}_{Gradient \ Coefficient}  \nabla_{\theta}\log \pi_{\theta}(o_t | q, o_{<t})\right).
\label{eq:objective}
\end{equation}

This framework is structured around three fundamental elements: 1) \textit{Data Source $\mathcal{D}$}, which specifies the set of data used for training; 2) \textit{Reward Function $\pi_{{rf}}$}, providing the feedback signal for learning; 3) \textit{Algorithm $\mathcal{A}$}, which transforms both the input data and reward signals into a gradient coefficient $GC$—influencing the intensity of updates, either as penalties or rewards, for the training examples. Building on this unified formulation, we summarize several typical approaches:

\begin{itemize}
    \item  \textbf{Supervised Fine-tuning (SFT)}: In this setting, the pretrained model is further trained on a dataset specifically curated by human annotators.
    \item \textbf{Rejection Sampling Fine-tuning (RFT)}: RFT continues the fine-tuning of the SFT model, but restricts the training data to those SFT model outputs that pass correctness filtering according to their responses to SFT prompts.
    \item \textbf{Direct Preference Optimization (DPO)}: DPO incrementally enhances the SFT model, leveraging pairwise comparison losses on additional outputs generated from the SFT model.
    \item \textbf{Online Rejection Sampling Fine-tuning (Online RFT)}: Unlike standard RFT, Online RFT begins with the SFT model as an initial policy and dynamically refines it using outputs sampled from the live (current) policy, rather than from the fixed SFT model.
    \item \textbf{PPO/GRPO}: Here, the policy model is initialized using the SFT parameters and further optimized via reinforcement, utilizing outputs sampled from the continuously updated policy.
\end{itemize}

We present a summary of the elements of these approaches in Table \ref{tab:data_policy}. For a comprehensive derivation, consult Appendix \ref{app:analysis-rl}.

\paragraph{Insight on Data Source} The data sources can be classified into two types: online sampling and offline sampling. Online sampling refers to scenarios in which training samples are drawn from the ongoing exploration by the current policy model, whereas offline sampling refers to using data collected from the outputs of the initial SFT model. Methods such as RFT and DPO utilize the offline approach, while Online RFT and GRPO adopt online sampling strategies.

Figure \ref{fig:rl-analysis} illustrates our findings that Online RFT yields substantially better performance compared to RFT across two benchmarks. Notably, during the initial training stages, Online RFT matches the performance of RFT, but as training progresses, it secures a clear performance lead, highlighting the benefits of online sampling. This pattern makes sense: early on, the policy (actor) and the SFT model remain quite similar, leading to minor variations in the collected samples. As training continues, however, the policy’s generated data diverges more notably from the SFT outputs, making the advantages of dynamic, real-time sampling increasingly pronounced.

\paragraph{Insight on Gradient Coefficient} The reward signal is transformed into a gradient coefficient, which guides updates to the model parameters. In our study, we distinguish the reward function into two categories: `Rule' and `Model'. A Rule-based reward evaluates response quality by directly checking the answer’s correctness, whereas the Model-based reward employs a learned reward model to assign a score to each response, using training data labeled via the rule-based assessment. Equations \ref{eq:GC-RFT} and \ref{eq:GC-GRPO} underscore a crucial distinction: GRPO uniquely modulates its gradient coefficient based on the magnitude of the reward from the reward model, enabling graded reinforcement and penalization depending on response quality. Conversely, Online RFT lacks this mechanism—incorrect answers are not penalized, and all correct responses are equally reinforced, regardless of their merit.

As illustrated in Figure \ref{fig:rl-analysis}, GRPO outperforms online RFT, underscoring the advantages of adjusting the coefficients for positive and negative gradients. Moreover, GRPO+PS achieves better outcomes than GRPO+OS, suggesting that leveraging fine-grained, step-aware gradient coefficients is advantageous. Additionally, we examine iterative RL in our experiments by applying two rounds of training. Figure \ref{fig:iter-rl} demonstrates that iterative RL leads to noticeable improvements in performance, particularly during the initial iteration.

\subsubsection{Why RL Works?} This work applies reinforcement learning on a subset of the instruction tuning data, resulting in substantial improvements over the baseline instruction tuning model. To gain further insight into the efficacy of reinforcement learning, we assess Pass@K and Maj@K accuracy for both Instruct and RL models using two benchmarks. Figure \ref{fig:maj-pass} reveals that RL increases Maj@K accuracy, yet does not affect Pass@K. This pattern suggests that RL strengthens the model’s output distribution by increasing the likelihood that correct responses appear within the TopK, rather than fundamentally enhancing core abilities. In a similar vein, \citep{wang2023making} identify a \textbf{misalignment problem} with reasoning performance in SFT models, showing that preference alignment methods can improve SFT’s reasoning abilities \citep{yuan2023rrhf,song2023preference,wang2023making}.

\subsubsection{How to Achieve More Effective RL?}
\label{sec:effec_rl}
Our results show that RL delivers strong results in mathematical reasoning scenarios. We also introduce a unified framework that encompasses a range of training strategies. Under this framework, various training methods can be seen as direct or approximate implementations of RL. As presented in Equation \ref{eq:objective}, three main aspects are emphasized: Data Source, Algorithm, and Reward Function. We outline several prospective research avenues regarding these elements.

\paragraph{Data Source} The data source constitutes the foundational material for all training techniques. For RL, this typically refers to unlabeled questions accompanied by outputs sampled from the current policy. In this study, we restrict ourselves to questions from the instruction tuning phase and use basic nucleus sampling to produce outputs. We hypothesize this may explain why improvements are seen chiefly in Maj@K performance. Moving forward, we plan to test the RL pipeline on prompts that differ from the original distribution and to adopt \textbf{advanced sampling (decoding) strategies}—such as tree-search-based techniques \citep{yao2023tree}. Furthermore, \textbf{efficient inference techniques} \citep{xia-etal-2023-speculative,leviathan2023fast,kwon2023efficient,xia2024unlocking}, which impact how efficiently the policy can explore, are expected to be crucial as well.

\paragraph{Algorithms} Algorithms utilize both the available data and the reward signal to calculate gradient coefficients that guide the adjustment of model parameters. According to Equation \ref{eq:objective}, contemporary methods fundamentally \textbf{TRUST} the feedback from the reward function to alter the conditional probability distribution over tokens. Nevertheless, the reliability of the reward signal cannot be guaranteed at all times, particularly for highly intricate tasks. For instance, the PRM800K datasets \citep{lightman2023let}—despite meticulous labeling efforts by experienced annotators—still contain around 20\% annotation errors\footnote{\url{https://github.com/openai/prm800k/issues/12\#issuecomment-1728491852}}. Consequently, it is imperative to investigate reinforcement learning algorithms that maintain robustness when exposed to noisy reward feedback. We anticipate that \textbf{WEAK-TO-STRONG} alignment strategies \citep{burns2023weak} have the potential to substantially reshape the nature of learning algorithms.

\paragraph{Reward Function} The reward function constitutes the main origin of the learning signal. Within reinforcement learning, this function is often embodied by a neural reward model. We identify three pivotal areas for further progress in reward modeling: 1) \textbf{Enhancing the reward model’s generalization capacity.} Effective reward models must generalize successfully to novel, out-of-distribution problems and advanced model outputs; absent this capability, reinforcement learning is liable to simply reinforce the existing distribution of large language models, without genuine improvements in their core competencies; 2) \textbf{Capturing the reward model’s uncertainty.} Modeling uncertainty could be instrumental as an intermediary between less reliable reward models and weak-to-strong algorithmic approaches; 3) \textbf{Constructing high-quality process reward models efficiently.} Such models are essential to supply granular training signals, especially for guiding complex reasoning steps \citep{lightman2023let,wang2023math}.

\section{Conclusion, Limitation, and Future Work}

In this work, we introduce DeepSeekMath, an open-source model that surpasses all other available open models on the MATH benchmark at the competition level and rivals the capabilities of proprietary systems. DeepSeekMath~builds upon DeepSeek-Coder-v1.5 7B as its base, followed by continued training spanning 500B tokens, with a substantial fraction (120B) consisting of mathematics-related data extracted from Common Crawl. Our comprehensive ablation analysis reveals that web content is a valuable reservoir for superior mathematical datasets, whereas data from arXiv contributes less than anticipated. We propose Group Relative Policy Optimization (GRPO), an alternative to Proximal Policy Optimization (PPO), which enhances mathematical reasoning while reducing memory requirements. Empirical evaluations demonstrate GRPO’s effectiveness, even at advanced performance levels of DeepSeekMath-Instruct 7B on standard benchmarks. Furthermore, we offer a unified framework for contextualizing a variety of reinforcement learning techniques, and we outline several promising avenues for advancing reinforcement learning methodologies.

Despite its notable achievements in quantitative reasoning tasks, DeepSeekMath~still lags behind closed models when addressing geometry and formal proof tasks. As observed in preliminary experiments, the model exhibits difficulty with geometric problems involving shapes such as triangles and ellipses, possibly reflecting a pre-training and fine-tuning data bias. Moreover, due to model size constraints, DeepSeekMath~does not match GPT-4’s capabilities in few-shot learning: while GPT-4 displays marked improvement with few-shot examples, DeepSeekMath~shows little difference between zero-shot and few-shot scenarios. Moving forward, we plan to refine our data selection strategy to compile a more robust and high-quality pre-training dataset. We also aim to investigate the promising research directions highlighted in Section \ref{sec:effec_rl}, with the goal of enhancing reinforcement learning strategies for large language models.